%----------
%	MODELOS EVALUADOS (multi-proveedor)
%----------

\label{sec:genai_models_iris}

La selección cubre las dos vías de despliegue y las principales familias
comerciales. En cada caso se eligen variantes ligeras y económicas, aptas para
procesar un corpus periodístico completo. La Tabla~\ref{tab:iris-modelos} recoge
los modelos considerados.

\begin{table}[H]
    \centering
    \begingroup
    \small
    \ttabbox[\FBwidth]{
        \caption{Modelos considerados y vía de acceso}
        \label{tab:iris-modelos}
    }{
        \begin{tabularx}{\textwidth}{
            l
            l
            >{\raggedright\arraybackslash}X
        }
            \hline
            \textbf{Proveedor} & \textbf{Modelo} & \textbf{Vía / observaciones} \\
            \hline
            Local (Ollama) & \texttt{gemma4:e4b} & Granja GPU TSC; modelo abierto de Google \\
            OpenAI & \texttt{gpt-4o-mini} & API \\
            OpenAI & \texttt{gpt-5.4-nano} & API; variante de razonamiento (desactivado) \\
            Google & \texttt{gemini-3.1-flash-lite} & API \\
            \hline
        \end{tabularx}
    }
    \endgroup
\end{table}

La evaluación del Capítulo~\ref{results_and_discussion} cubre estos cuatro
modelos: \texttt{gemma4:e4b} en local y las tres variantes comerciales de API.
Tres decisiones sobre la selección merecen justificarse.

Para que la comparación entre proveedores sea justa, se desactivó el razonamiento
explícito en los modelos que lo admiten (el parámetro \texttt{reasoning\_effort}
en OpenAI, \texttt{thinking\_budget} en Google). Así ningún modelo dispone de
cómputo extra del que carezcan los demás.

Los modelos locales no cuestan dinero, porque corren en infraestructura propia,
pero sí tienen un coste computacional alto que condiciona su viabilidad a escala
(Sección~\ref{subsec:iris_coste}).

La preferencia por variantes ligeras (\textit{mini}, \textit{nano},
\textit{flash-lite}, modelos locales de pocos miles de millones de parámetros) es
deliberada. Un sistema de asistencia editorial que procese miles de piezas
necesita un coste por documento contenido. Esa elección tiene un precio, menor
capacidad frente a los modelos punteros, que se discute a la luz de los
resultados en el Capítulo~\ref{results_and_discussion}.
